\scanpage{13}\mspage{12\unsure{..|.}}
% ink top, first two lines
\unsure{et|Et} choisissons une ``origine'' $a$ dans $\unsure{E|\mathcal E}=\{a,b\}$ (ce qui va déterminer la relation d'ordre).

% ink upper third, Considérons
\deleted{Soit} \deleted{Alors} Considérons sur $\mathcal L$ la relation d'ordre \unsure{définissant} le biordre, telle que $a<b$. On va voir qu'elle est \unsure{entièmt|entièrement} \unsure{déterminée} par la conn. \deleted{\unsure{de}}\unsure{du} \deleted{\unsure{la paire}}\inserted{couple} $(a,b)$

% ink upper middle, On aura, three defining sets
On \unsure{aura|a} \unsure{.}
% slanted
\margin{déterminés par $R$}
\[
\left\{\begin{aligned}
\mathcal L_{\deleted{\unsure{>}}\geq b}&=\{\,x\in L\mid\deleted{x=b\text{ ou}}\ \inserted{x=b\text{ ou}}\ b\text{ entre }a\text{ et }x\,\}\quad\deleted{\unsure{a,b|a<b}.}\\
\mathcal L_{\leq a}&=\{\,x\in L\mid\inserted{x=a\text{ ou}}\ a\text{ entre }\deleted{\unsure{a}}b\text{ et }x\,\}\\
\deleted{\unsure{\mathcal L}}\mathopen{]}\deleted{\unsure{x}}a,b\mathclose{[}&=\{\,x\in L\mid x\text{ entre }a\text{ et }b\,\}
\end{aligned}\right.
\]
% ink left, below brace
% slanted
\margin{\unsure{donc}}

% ink middle, question
\deleted{Quand}\inserted{Soient} $u,v\in L$, quand a-t-on
\[
u\leq v\qquad\unsure{\mathord{?}|.}
\]
Il faut et suffit qu'il existe
\[
\begin{aligned}
&\deleted{x\in}\quad x,y\in L\quad\text{tels que}\\
&x\in L_{\leq a},\ y\in L_{\geq b}\qquad\text{et}\\
&\hspace{2em}x\leq u\leq v\leq y
\end{aligned}
\]
% ink bottom, i.e., left bracket
\[
\text{i.e.}\quad
\left[\begin{aligned}
&\deleted{\{\unsure{x},y,\unsure{u|b},v\}\ \text{est }\in\mathrm{Drap}(\mathcal L)}\\
&\qquad\deleted{x,y\in}\quad x\in L_{\leq a},\ y\in L_{\geq b}\quad(\text{donc }x\unsure{<|\leq}y)\\
&\deleted{\text{avec }x\text{ plus petit élément}}\\
&\deleted{u\text{ est entre}}\quad u,v\text{ sont entre }x\text{ et }y\\
&\quad u\text{ est entre }x\text{ et }v\qquad\unsure{.}\\
&\quad v\text{ est entre }u\text{ et }y
\end{aligned}\right.
\]
